\documentclass[11pt]{article}

\usepackage[utf8]{inputenc}
\usepackage[T1]{fontenc}
\usepackage{lmodern}
\usepackage[a4paper,margin=1in]{geometry}
\usepackage{amsmath,amssymb,amsthm}
\usepackage{booktabs}
\usepackage{graphicx}
\usepackage{hyperref}
\usepackage[nameinlink,noabbrev]{cleveref}
\usepackage{xcolor}

\hypersetup{
  colorlinks=true,
  linkcolor=blue!50!black,
  citecolor=blue!50!black,
  urlcolor=blue!50!black
}

\title{Zero-Knowledge Proofs for Gradient Boosted Decision Trees}
\author{Author Name}
\date{}

\numberwithin{equation}{section}
\newtheorem{theorem}{Theorem}[section]
\newtheorem{lemma}[theorem]{Lemma}
\newtheorem{proposition}[theorem]{Proposition}
\newtheorem{corollary}[theorem]{Corollary}
\theoremstyle{definition}
\newtheorem{definition}[theorem]{Definition}
\theoremstyle{remark}
\newtheorem{remark}[theorem]{Remark}

\begin{document}

\maketitle

\begin{abstract}
We study zero-knowledge proofs for quantized gradient boosted decision trees
(GBDTs). The main bottleneck is not tree inference itself, but proving that each
boosting round uses the split that maximizes the training objective without
reducing the entire algorithm to a large generic circuit. Our approach keeps the
tree-specific structure explicit. We use a common batched zero-check template to
combine polynomial identities across heterogeneous subgroup sizes, together with
replication and interleaving operators for reshaping committed tensors. On top
of this compatibility layer, we introduce aggregation and histogram relations
that expose split search as algebra over per-node, per-feature, and per-bin
statistics. This leads to a training proof organized around a small catalog of
relations: lookup, tensor sums, aggregation, histogram construction, prefix
sums, zero-safe division, and max selection. The resulting architecture is
aimed at security-focused systems venues: it preserves the structure of GBDT
training, cleanly separates imported subroutines from new primitives, and
supports both batched inference and end-to-end training proofs.
\end{abstract}

\section{Introduction}

Zero-knowledge proofs for machine learning systems are increasingly relevant
when model correctness, privacy, and auditability must coexist. This section
 can introduce the motivation for proving correct GBDT inference or training,
the limitations of prior approaches, and the high-level idea behind your
construction.

\paragraph{Suggested questions to answer.}
Why is ZK for GBDT important? What is technically difficult? What does your
approach achieve that prior work does not?

\section{Preliminaries}

\subsection{Notation and Polynomial Encoding}

We work over a finite field $\mathbb{F}$. All vectors, matrices, and tensors are
flattened in row-major order before commitment. For $n \in \mathbb{N}$, we write
$[n] := \{0,1,\ldots,n-1\}$.

All committed lengths are assumed to be powers of two. To make this precise, we
fix an integer $L \ge 1$ such that every committed object is smaller than $2^L$,
and is padded to a length
dividing $2^L$. We also fix a primitive $2^L$-th root of unity $\omega \in
\mathbb{F}$. For every power of two $n$ with $n \mid 2^L$, we define
\[
\omega_n := \omega^{2^L/n}
\qquad\text{and}\qquad
\mathbb{H}_n := \langle \omega_n \rangle.
\]
Thus $\mathbb{H}_n$ is the multiplicative subgroup of size $n$ used as the
evaluation domain for length-$n$ vectors. Its vanishing polynomial is
\[
\nu_n(X) := \prod_{a \in \mathbb{H}_n}(X-a) = X^n-1.
\]
Because all domains are derived from the same root $\omega$, we have
$\omega_{nm}^m=\omega_n$ whenever $n$, $m$, and $nm$ are powers of two dividing
$2^L$. This compatibility is used later when we align objects of different
shapes.

Let $\lambda_n^{(i)}(X)$ denote the $i$-th Lagrange basis polynomial on
$\mathbb{H}_n$.

For a vector $\mathbf{v}\in\mathbb{F}^n$, its polynomial encoding over
$\mathbb{H}_n$ is
\[
v(X) := \sum_{i \in [n]} \mathbf{v}[i] \lambda_n^{(i)}(X),
\]
where $\mathbf{v}[i]$ denotes the $i$-th entry of $\mathbf{v}$.

Throughout the paper, we need to manipulate the shapes
of a high-dimensional tensor. Given $\mathbf{V}\in \mathbb{F}^{a\times b\times c}$,
we write $\mathbf{V}_{\{abc\}}\in \mathbb{F}^{abc}$ for its flatten form,
and $\mathbf{V}_{\{a\times bc\}}\in \mathbb{F}^{a\times bc}$ for its flatten form with the first dimension preserved.
As a sanity check, note that
\[
\mathbf{V}_{\{a\times b\times c\}}[i, j, k] = \mathbf{V}_{\{a\times bc\}}[i, j \cdot c + k] = \mathbf{V}_{\{abc\}}[i \cdot bc + j \cdot c + k].
\]
A similar notation applies to tensors of other dimensions.

\begin{table*}[t]
\centering
\small
\begin{tabular}{lllll}
\toprule
Symbol & Shape & Support & Meaning & Status \\
\midrule
$\mathbf{D}$ & $M\times d$ & all samples & quantized dataset & committed \\
$\mathbf{Y}$ & $M\times K$ & all samples & one-hot labels & committed \\
$\mathbf{S}^{(t)}$ & $M\times K$ & round $t$ & cumulative scores & committed slice \\
$\mathbf{G}^{(t)}$ & $M\times K$ & round $t$ & pseudo-residuals & committed slice \\
$\boldsymbol{\ell}^{(t,k)}$ & $M$ & tree $(t,k)$ & leaf assignment & committed slice \\
$\mathbf{C}^{(t,k)},\mathbf{H}^{(t,k)}$ & $N_{\mathrm{tot}}\times d\times B$ & tree $(t,k)$ & count / sum histograms & committed slice \\
$\Gamma^{(t,k)}$ & $N_{\mathrm{int}}\times d\times B$ & tree $(t,k)$ & candidate split scores & committed slice \\
$\widetilde{\mathbf{W}}^{(t)}$ & $N_{\mathrm{tot}}\times K$ & round $t$ & padded leaf values & committed slice \\
$F(X),\widehat{u}(X)$ & varies & derived & batched virtual polynomials & virtual \\
\bottomrule
\end{tabular}
\caption{Main proof-facing objects used later in Sections~3--5. All tensors are
flattened in row-major order before commitment.}
\label{tab:object-glossary}
\end{table*}

\subsection{KZG Polynomial Commitments}

We use a KZG-based polynomial commitment scheme (PCS) for univariate polynomials.
The PCS consists of the following algorithms:

\begin{itemize}
    \item $\Setup(1^\lambda, d_{\max}) \rightarrow \pp$:
    Generates public parameters supporting commitments to polynomials of degree at most $d_{\max}$.

    \item $\Commit(\pp, f) \rightarrow c$:
    Commits to a polynomial $f$ and outputs a commitment $c$.

    \item $\Open(\pp, f, z) \rightarrow (v, \pi)$:
    Given a polynomial $f$ and an evaluation point $z$, outputs the evaluation
    $v = f(z)$ together with an opening proof $\pi$.

    \item $\Verify(\pp, c, z, v, \pi) \rightarrow \{0,1\}$:
    Verifies that $v$ is the correct evaluation of the polynomial committed in $c$
    at point $z$.
\end{itemize}

\paragraph{Batched openings.}
In addition to single-point openings, we use a batched opening interface that
allows the prover to open multiple evaluation claims together.

\begin{itemize}
    \item $\BatchOpen(\pp, \{f_i\}_{i \in [m]}, \{Q_i\}_{i \in [m]})
    \rightarrow (\{v_{i,z}\}_{i \in [m], z \in Q_i}, \Pi)$:
    For each committed polynomial $f_i$ and each query point set $Q_i$,
    outputs all claimed evaluations $v_{i,z} = f_i(z)$ together with a single
    batched proof $\Pi$.

    \item $\BatchVerify(\pp, \{c_i\}_{i \in [m]}, \{Q_i\}_{i \in [m]},
    \{v_{i,z}\}_{i \in [m], z \in Q_i}, \Pi) \rightarrow \{0,1\}$:
    Verifies that all claimed evaluations are consistent with the corresponding
    commitments.
\end{itemize}

\paragraph{Properties.}
We rely on the following standard properties. \emph{Completeness}: honestly generated openings
and batched openings always verify.  \emph{Binding}: after a
commitment is published, the prover cannot later open it to a different value at
the same point. 

\subsection{Gradient Boosted Decision Trees in Plaintext}

Given a commitment to a forest, the task of proof of training is to show that
the forest arises from a valid execution of the GBDT training procedure.

Let $\mathbf{D}\in[B]^{M\times d}$ be the quantized training dataset, where
$M$ is the number of samples, $d$ is the number of features, and each feature
value lies in one of $B$ bins. Let $\mathbf{Y}\in\{0,1\}^{M\times K}$ be the
one-hot label matrix for a $K$-class classification task. We treat the number
of boosting rounds $T$,
the number of classes $K$,
and the size of dataset $M$
as public parameters of the statement being proved.

\paragraph{Plaintext training.}
GBDT trains an additive model over $T$ rounds. In each round $t\in[T]$, the
learner maintains a score matrix $\mathbf{S}^{(t)}\in\mathbb{F}^{M\times K}$,
where $\mathbf{S}^{(t)}[r,k]$ denotes the current score of sample $r$ for class
$k$, before adding the trees of round $t$. We initialize
$\mathbf{S}^{(0)}=\mathbf{0}$.

The current class probabilities are obtained by applying softmax row-wise:
\begin{equation}
\boldsymbol{\pi}^{(t)}[r,k]
:=
\frac{\exp(\mathbf{S}^{(t)}[r,k])}
{\sum_{k'\in[K]} \exp(\mathbf{S}^{(t)}[r,k'])}.
\end{equation}
Let $\mathbf{Y}\in\{0,1\}^{M\times K}$ be the one-hot label matrix. The
multinomial logistic loss at round $t$ is
\begin{equation}
\mathcal{L}^{(t)}
:=
-\sum_{r\in[M]}\sum_{k\in[K]}
\mathbf{Y}[r,k]\log \boldsymbol{\pi}^{(t)}[r,k].
\end{equation}
Its partial derivative with respect to the score entry
$\mathbf{S}^{(t)}[r,k]$ is
\begin{equation}
\frac{\partial \mathcal{L}^{(t)}}{\partial \mathbf{S}^{(t)}[r,k]}
=
\boldsymbol{\pi}^{(t)}[r,k]-\mathbf{Y}[r,k].
\end{equation}
Therefore, the negative gradient is
\begin{equation}\label{eq:plain-pseudoresidual}
\mathbf{G}^{(t)}[r,k]
:=
-\frac{\partial \mathcal{L}^{(t)}}{\partial \mathbf{S}^{(t)}[r,k]}
=
\mathbf{Y}[r,k]-\boldsymbol{\pi}^{(t)}[r,k].
\end{equation}
We call $\mathbf{G}^{(t)}$ the pseudo-residual matrix. Similar to the gradient descent
approach, we want to move the score matrix towards the negative gradient direction, and therefore reduce the loss.
Thus in round $t$, the
learner trains one regression tree for each class $k\in[K]$ to fit the target
vector $\mathbf{G}^{(t)}[\cdot,k]\in\mathbb{F}^{M}$.\footnote{The above describes the training process for typical gradient-boosting decision trees. 
It is also possible to adapt our constructions to other variants,
e.g. XGBoost~\cite{xgboost16},
by modifying the training target and tree-splitting policy accordingly.
See \textcolor{red}{Appendix}.}

\paragraph{Regression tree.}
Fix a round $t$ and a class $k$. The learner constructs a regression tree
$f^{(t,k)}:[B]^d\to\mathbb{F}$, which is a binary decision tree whose leaves
carry scalar prediction values. Each internal node $u$ is labeled by a feature
index $j_u\in[d]$ and a threshold bin $b_u\in[B]$. Given a sample
$\mathbf{x}\in[B]^d$, routing at node $u$ is determined by the predicate
$\mathbf{x}[j_u] < b_u$: if it holds, the sample is sent to the left child of
$u$; otherwise it is sent to the right child. The output
$f^{(t,k)}(\mathbf{x})$ is the value stored at the leaf reached by this routing
procedure.

Similarly, during training, each node $u$ is associated with the set
$I_u\subseteq[M]$ of samples currently reaching $u$. The root contains all
samples, i.e., $I_{\mathrm{root}}=[M]$, and each split partitions the parent
sample set into the left and right child sample sets.

\paragraph{How a node is split.}
For the class-$k$ tree in round $t$, let
\[
g_r := \mathbf{G}^{(t)}[r,k]
\]
denote the regression target attached to sample $r$. Consider a node $u$ with
current sample set $I_u$. For a candidate split specified by a feature
$j\in[d]$ and a threshold bin $b\in[B]$, define
\begin{align}
L_{u,j,b} &:= \{\, r\in I_u : \mathbf{D}[r,j] < b \,\},\\
R_{u,j,b} &:= \{\, r\in I_u : \mathbf{D}[r,j] \ge b \,\}.
\end{align}
The quality of this split is measured by the residual sum of squares (RSS):
\begin{equation}\label{eq:rss-plain}
\operatorname{RSS}(u,j,b)
=
\sum_{r\in L_{u,j,b}} (g_r-\bar{g}_L)^2
+
\sum_{r\in R_{u,j,b}} (g_r-\bar{g}_R)^2,
\end{equation}
where
\[
\bar{g}_L := \frac{1}{|L_{u,j,b}|}\sum_{r\in L_{u,j,b}} g_r,
\qquad
\bar{g}_R := \frac{1}{|R_{u,j,b}|}\sum_{r\in R_{u,j,b}} g_r
\]
are the empirical means on the two sides. Intuitively, the split is good if the
targets assigned to each child are well concentrated around a single mean.

Among all valid candidates, i.e. those for which both $L_{u,j,b}\neq\emptyset$ and $R_{u,j,b}\neq\emptyset$, the trainer selects a
feature-threshold pair minimizing \eqref{eq:rss-plain}.
If valid split is not possible, the node is not split and becomes a leaf.

Once the tree structure is fixed, each leaf $\ell$ is assigned the mean target
value of the samples that reach it. Writing $I_\ell\subseteq[M]$ for the sample
set of leaf $\ell$, its output value is
\begin{equation}\label{eq:plain-leaf-value}
w_{\ell}^{(t,k)}
:=
\frac{1}{|I_\ell|}\sum_{r\in I_\ell} \mathbf{G}^{(t)}[r,k].
\end{equation}
Therefore, for every training sample $r$, the tree prediction
$f^{(t,k)}(\mathbf{D}[r,\cdot])$ is exactly the leaf value associated with the
leaf reached by sample $r$.

\paragraph{Round update.}
After training all $K$ regression trees
$\{f^{(t,k)}\}_{k\in[K]}$ in round $t$, the learner updates the score matrix by
adding the tree outputs with learning rate $\mu>0$:
\begin{equation}\label{eq:plain-score-update}
\mathbf{S}^{(t+1)}[r,k]
=
\mathbf{S}^{(t)}[r,k] + \mu\, f^{(t,k)}(\mathbf{D}[r,\cdot]).
\end{equation}
This completes round $t$. The next round recomputes
$\boldsymbol{\pi}^{(t+1)}$, forms the new pseudo-residual matrix
$\mathbf{G}^{(t+1)}$ using \eqref{eq:plain-pseudoresidual}, and trains a fresh
collection of $K$ regression trees against those updated targets. After $T$
rounds, the forest consists of the $TK$ trees trained all rounds.

At inference time, for a fresh sample $\mathbf{x}\in[B]^d$, the model
accumulates the per-class scores
\[
\sum_{t=0}^{T-1} f^{(t,k)}(\mathbf{x})
\qquad\text{for each }k\in[K],
\]
and predicts from the resulting $K$-dimensional score vector, typically by
applying softmax.

\paragraph{Implication for the proof.}
The proof-critical step is the correctness of the regression-tree training at
each node, namely, that the committed split metadata and leaf values are
consistent with the pseudo-residual targets of that round and class. In
standard plaintext implementations, the optimizer often evaluates split
candidates by sorting samples according to each feature and scanning candidate
thresholds in that sorted order. In quantized systems, the same optimization is
commonly expressed through per-bin histograms. Our proof follows the latter
view: instead of proving a sorting procedure in zero knowledge, we prove the
correctness of the histogram statistics from which the trainer derives the same
split scores.

\subsection{Quantization}

We follow the fixed-point quantization approach used in prior ZKP for ML works~\cite{sparrow24,zkdt20,zkboost26}. 
Given a real number $x$, we encode it as a signed integer
\[
q_x = \lfloor x \cdot 2^Q \rfloor,
\]
where $Q>0$ is a public quantization parameter. The corresponding fixed-point value is then interpreted as $x \approx q_x / 2^Q$.

Under this representation, addition is scale-preserving: if
$x_1 \approx q_{x_1}/2^Q$ and $x_2 \approx q_{x_2}/2^Q$, then
\[
x_1 + x_2 \approx \frac{q_{x_1}+q_{x_2}}{2^Q}.
\]
Multiplication, however, increases the scale and therefore requires an explicit
rescaling step. In particular, if
$x \approx q_x/2^Q$ and $y \approx q_y/2^Q$, then their product is represented by
\[
q_z = \left\lfloor \frac{q_x q_y}{2^Q} \right\rfloor,
\]
so that
\[
xy \approx \frac{q_z}{2^Q}.
\]
\subsection{Statement and Proof Overview}

Our goal is to prove correct training, not merely well-formedness of the
committed forest. Concretely, given commitments to a training dataset and to a
forest, the prover must show that the committed forest is exactly the result of
running the public quantized GBDT training algorithm on the committed dataset.

The public statement consists of the training parameters
$(M,d,B,K,T)$, the learning rate $\mu$, the public tree-shape policy, the
quantization parameters, the public tables used by imported subprotocols, the
KZG setup parameters, and the commitments to the dataset and the forest. The
witness consists of the values needed to certify one valid training trace,
including the labels, the round-wise boosting states, the histogram
statistics, the split metadata, the leaf values, and the auxiliary witnesses
required by the imported gadgets.

Accordingly, the verifier checks the following claim: there exists a witness
such that, starting from the committed dataset, the public quantized training
procedure produces exactly the committed forest. Thus, an accepting proof
certifies not only that the committed forest has valid syntax, but that it is
consistent with one end-to-end execution of the trainer semantics fixed by this
paper.

At a high level, the proof of training must establish four facts:
\begin{enumerate}
    \item the committed forest derives the claimed inference over the dataset;
    \item each boosting round correctly derives the pseudo-residuals and updates
    the score state;
    \item in each boosting round, each tree is trained according to the public training algorithm.
\end{enumerate}

The rest of the protocol decomposes this global claim into local checks for
forest semantics, routing, histogram construction, split-score correctness,
split optimality, and leaf-value correctness, and then batches these checks
into a small number of polynomial openings.

\begin{theorem}[Informal security statement]
\label{thm:main-informal}
Assume the KZG commitment scheme is binding, and the imported subprotocols satisfy
their claimed completeness and soundness properties. Then the protocol satisfies:
\begin{enumerate}
    \item \textbf{Completeness.} If the committed forest is obtained by honestly
    running the public quantized trainer on the committed dataset, then the
    prover can produce an accepting proof.

    \item \textbf{Soundness.} If the verifier accepts, then, except with the
    accumulated soundness error of the imported gadgets and the batching
    checks, the committed forest is consistent with a valid execution of that
    same trainer on the committed dataset.

    \item \textbf{(Optional) Witness privacy.} If the hiding variant of the commitment
    scheme is used and the imported subprotocols are zero knowledge, then the
    proof reveals nothing beyond the public statement and the committed
    outputs.
\end{enumerate}
We claim soundness, not a separate argument-of-knowledge statement.
\end{theorem}

\section{Technical Overview}

The proof system is best viewed as a certification procedure for quantized GBDT
training. The main challenge is not tree inference itself, but proving that
every boosting round computes the correct split statistics and commits a split
that optimizes the training objective. Two obstacles recur throughout the
construction. First, the proof naturally produces claims over different
polynomial domains and tensor layouts. Second, plaintext split search is
usually implemented through sorting or feature scans, which are awkward proof
targets. The overview below therefore isolates three ingredients: the
committed forest interface, the algebraic rewriting of split search, and the
batching tools that make the resulting checks compose succinctly.

\subsection{Forest Representation and Batched Inference}

We adopt the hyperbox forest encoding of \cite{cq++24}. Each node carries a
feature-wise half-open interval box, and a sample reaches a leaf exactly when
its quantized feature vector lies inside that box. All committed objects share
one global node namespace, so the row index
$u \in [N_{\mathrm{tot}}]$ refers to the same node everywhere. The committed
forest consists of:
\begin{itemize}
  \item routing matrices
  $\mathbf{L},\mathbf{R}\in\{0,1\}^{N_{\mathrm{tot}}\times N_{\mathrm{tot}}}$,
  where an internal-node row is one-hot and selects the left or
  right child of that node, while leaf rows are zero;
  \item lower and upper node bounds
  $\mathbf{N}^{\leftarrow},\mathbf{N}^{\rightarrow}\in[B+1]^{N_{\mathrm{tot}}\times d}$,
  where row $u$ stores the $d$-dimensional half-open box of node $u$ and every
  leaf row describes one accepting region of the tree;
  \item a feature-selection matrix
  $\mathbf{E}\in\{0,1\}^{N_{\mathrm{tot}}\times d}$,
  whose internal-node rows are one-hot and indicate which feature is split at
  each node, while leaf rows are zero;
  \item a threshold-index vector
  $\boldsymbol{\Theta}\in[B]^{N_{\mathrm{tot}}}$,
  whose internal-node entries record the selected threshold bin and whose leaf
  and padding rows are zero;
  \item a tree-index vector $\mathbf{tree}\in[TK]^{N_{\mathrm{tot}}}$ and a
  root matrix $\mathbf{Root}\in\{0,1\}^{TK\times N_{\mathrm{tot}}}$;
  \item a leaf-score vector $\mathbf{score}\in\mathbb{F}^{N_{\mathrm{tot}}}$,
  whose leaf entries store the scalar contribution returned by that tree and
  whose internal-node entries are zero or padding.
\end{itemize}
The pair $(\mathbf{E},\boldsymbol{\Theta})$ witnesses that
$\mathbf{N}^{\leftarrow}$ and $\mathbf{N}^{\rightarrow}$ encode a valid
hyperbox forest. The routing matrices $(\mathbf{L},\mathbf{R})$ are reused in
training to propagate histogram statistics from leaves to internal nodes.

\Cref{fig:forest-encoding-inference} follows the exposition style of
\cite{cq++24}: the left panel presents two concrete trees in the global node
namespace, while the right panel shows the committed row view on which the
proof operates. The committed columns record the tree index, split metadata,
node box, and leaf score. The online fetch opens from this view only the tree
index, node box, and score of the accepting leaf.

\begin{figure*}[t]
  \centering
  \resizebox{\linewidth}{!}{%
  \begin{tikzpicture}[
    >=Latex,
    font=\small,
    paneltitle/.style={font=\bfseries},
    treelabel/.style={font=\bfseries\small},
    internal/.style={
      circle,
      draw=blue!55!black,
      line width=0.55pt,
      fill=blue!5,
      minimum size=10.2mm,
      inner sep=1pt,
      align=center
    },
    leaf/.style={
      rounded corners=2pt,
      draw=green!50!black,
      line width=0.55pt,
      fill=green!10,
      minimum width=14mm,
      minimum height=8.6mm,
      inner sep=2pt,
      align=center
    },
    note/.style={align=center, font=\scriptsize},
    bigtable/.style={
      matrix of nodes,
      ampersand replacement=\&,
      row sep=0.75mm,
      column sep=\colgap,
      nodes={
        anchor=center,
        inner xsep=2.5pt,
        inner ysep=1.2pt,
        minimum height=4.5mm,
        text depth=.25ex,
        outer sep=0pt
      },
      row 1/.style={nodes={font=\bfseries}},
      column 1/.style={nodes={minimum width=7mm}},
      column 2/.style={nodes={minimum width=11mm}},
      column 3/.style={nodes={minimum width=10mm}},
      column 4/.style={nodes={minimum width=13mm}},
      column 5/.style={nodes={minimum width=9mm}},
      column 6/.style={nodes={minimum width=13mm}},
      column 7/.style={nodes={minimum width=13mm}},
      column 8/.style={nodes={minimum width=12mm}}
    },
    colbg/.style={
      draw=gray!70,
      rounded corners=2pt,
      line width=.5pt,
      fill=white,
      inner sep=2.8pt
    },
    meta/.style={
      draw=gray!70,
      rounded corners=2pt,
      fill=gray!6,
      align=center,
      inner sep=6pt
    },
    fetchrow/.style={
      draw=red!75!black,
      rounded corners=2pt,
      line width=.85pt,
      inner xsep=4pt,
      inner ysep=2pt
    },
    fetchtxt/.style={text=red!75!black, font=\bfseries},
    outer/.style={
      draw=gray!55,
      dashed,
      rounded corners=4pt,
      inner sep=8pt
    }
  ]

    \def\colgap{4.2mm}

    % =========================
    % Left panel: tree semantics
    % =========================
    \node[paneltitle] (ltitle) at (-4.8,3.12) {(a) Tree semantics};

    % Tree 1
    \node[treelabel] at (-7.55,2.42) {tree $1$};
    \node[internal] (t0r) at (-7.55,1.72) {$1$\\[-1mm]\scriptsize$x_1<3$};
    \node[internal] (t0l) at (-9.05,0.38) {$2$\\[-1mm]\scriptsize$x_2<4$};
    \node[leaf]     (t0rleaf) at (-6.00,0.38) {$7$\\[-1mm]\scriptsize$w_2$};
    \node[leaf]     (t0w0) at (-10.00,-1.02) {$5$\\[-1mm]\scriptsize$w_0$};
    \node[leaf]     (t0w1) at (-8.35,-1.02) {$6$\\[-1mm]\scriptsize$w_1$};

    \draw[->, line width=.75pt] (t0r) -- (t0l);
    \draw[->, line width=.75pt] (t0r) -- (t0rleaf);
    \draw[->, line width=.75pt] (t0l) -- (t0w0);
    \draw[->, line width=.75pt] (t0l) -- (t0w1);

    % Tree 2
    \node[treelabel] at (-2.65,2.42) {tree $2$};
    \node[internal] (t1r) at (-2.65,1.72) {$3$\\[-1mm]\scriptsize$x_1<2$};
    \node[leaf]     (t1lleaf) at (-4.15,0.38) {$8$\\[-1mm]\scriptsize$w'_0$};
    \node[internal] (t1rr) at (-1.15,0.38) {$4$\\[-1mm]\scriptsize$x_2<6$};
    \node[leaf]     (t1w1) at (-2.05,-1.02) {$9$\\[-1mm]\scriptsize$w'_1$};
    \node[leaf]     (t1w2) at (-0.40,-1.02) {$10$\\[-1mm]\scriptsize$w'_2$};

    \draw[->, line width=.75pt] (t1r) -- (t1lleaf);
    \draw[->, line width=.75pt] (t1r) -- (t1rr);
    \draw[->, line width=.75pt] (t1rr) -- (t1w1);
    \draw[->, line width=.75pt] (t1rr) -- (t1w2);

    % =====================================
    % Right panel: one unified committed table
    % =====================================
    \node[paneltitle] (rtitle) at (7.25,3.12) {(b) Committed row view used by the proof};

    \matrix[bigtable, anchor=north west] (tbl) at (1.15,2.55) {
      $\mathbf{u}$ \& $\mathbf{tree}$ \& $\mathbf{is\_leaf}$ \& $\mathbf{E}$ \& $\boldsymbol{\Theta}$ \& $\mathbf{N}^{\leftarrow}$ \& $\mathbf{N}^{\rightarrow}$ \& $\mathbf{score}$ \\
      $0$ \& $0$ \& $0$ \& $(0,0)$ \& $0$ \& $(0,0)$ \& $(0,0)$ \& $0$ \\
      $1$ \& $1$ \& $0$ \& $(1,0)$ \& $3$ \& $(0,0)$ \& $(B,B)$ \& $0$ \\
      $2$ \& $1$ \& $0$ \& $(0,1)$ \& $4$ \& $(0,0)$ \& $(3,B)$ \& $0$ \\
      $3$ \& $2$ \& $0$ \& $(1,0)$ \& $2$ \& $(0,0)$ \& $(B,B)$ \& $0$ \\
      $4$ \& $2$ \& $0$ \& $(0,1)$ \& $6$ \& $(2,0)$ \& $(B,B)$ \& $0$ \\
      $5$ \& $1$ \& $1$ \& $(0,0)$ \& $0$ \& $(0,0)$ \& $(3,4)$ \& $w_0$ \\
      $6$ \& $1$ \& $1$ \& $(0,0)$ \& $0$ \& $(0,4)$ \& $(3,B)$ \& $w_1$ \\
      |[fetchtxt]| $7$ \& |[fetchtxt]| $1$ \& |[fetchtxt]| $1$ \& |[fetchtxt]| $(0,0)$ \& |[fetchtxt]| $0$ \& |[fetchtxt]| $(3,0)$ \& |[fetchtxt]| $(B,B)$ \& |[fetchtxt]| $w_2$ \\
      $8$ \& $2$ \& $1$ \& $(0,0)$ \& $0$ \& $(0,0)$ \& $(2,B)$ \& $w'_0$ \\
      |[fetchtxt]| $9$ \& |[fetchtxt]| $2$ \& |[fetchtxt]| $1$ \& |[fetchtxt]| $(0,0)$ \& |[fetchtxt]| $0$
        \& |[fetchtxt]| $(2,0)$ \& |[fetchtxt]| $(B,6)$ \& |[fetchtxt]| $w'_1$ \\
      $10$ \& $2$ \& $1$ \& $(0,0)$ \& $0$ \& $(2,6)$ \& $(B,B)$ \& $w'_2$ \\
    };

    % Column background boxes
    \begin{scope}[on background layer]
      \foreach \c in {1,...,8}{
        \node[colbg, fit=(tbl-1-\c)(tbl-12-\c)] {};
      }
    \end{scope}

    % Horizontal separators: after header and after internal rows
    \foreach \r in {1,6}{
      \draw[gray!60, line width=.4pt]
        ([xshift=-2pt]tbl-\r-1.south west) -- ([xshift=2pt]tbl-\r-8.south east);
    }

    % Highlight fetched rows (one row per tree)
    \node[fetchrow, fit=(tbl-9-1)(tbl-9-8)] (fetchA) {};
    \node[fetchrow, fit=(tbl-11-1)(tbl-11-8)] (fetchB) {};

    % Query box
    \node[meta, anchor=north west, text width=3.35cm] (sample)
      at ([xshift=6mm]tbl.north east)
      {query $\mathbf{x}=(4,2)$\\[1mm]
       fetch one row per tree:\\
        $u_1=7,\;u_2=9$\\[1mm]
       and prove, for any fetched row $u$,\\
       $\mathbf{N}^{\leftarrow}[u]\le \mathbf{x}<\mathbf{N}^{\rightarrow}[u]$};

    % Arrows to the two fetched rows
    \draw[->, very thick, red!75!black]
      ([yshift=7pt]sample.west) to[out=180,in=0] ([xshift=2pt]fetchA.east);
    \draw[->, very thick, red!75!black]
      ([yshift=-7pt]sample.west) to[out=180,in=0] ([xshift=2pt]fetchB.east);

    % Outer frame
    \node[
      outer,
      fit=(ltitle)(rtitle)
          (t0r)(t0l)(t0rleaf)(t0w0)(t0w1)
          (t1r)(t1lleaf)(t1rr)(t1w1)(t1w2)
          (tbl)(sample)
    ] {};

  \end{tikzpicture}%
  }
  \caption{An illustration of the forest encoding and inference semantics.
  The left panel shows two compact trees with
  global node identifiers, and the right panel shows the committed
  objects used by inference. The highlighted entries are the fetched rows for one
  query, one per tree. Batched inference performs this fetch for every
  tree-sample pair in one shot and then aggregates scores across trees for
  each sample-class pair.}
  \Description{Two example decision trees and their committed table
  representation. The table includes node identifiers, tree identifiers,
  leaf indicators, split-feature masks, thresholds, hyper-box bounds, and leaf
  scores. Two leaf rows are highlighted to show how inference fetches one row
  per tree and checks interval containment for a query point.}
  \label{fig:forest-encoding-inference}
\end{figure*}


This encoding supports both inference and training. For inference, the protocol
follows the row-fetch pattern of \cite{cq++24}. Given a query batch
$\mathbf{X}\in[B]^{b\times d}$, the prover conceptually replicates
$\mathbf{X}$ once per tree, determines the accepting leaf row of each
tree-sample pair, and proves a batched fetch of
\[
(\mathbf{tree},\mathbf{N}^{\leftarrow},\mathbf{N}^{\rightarrow},\mathbf{score})
\]
from those committed leaf rows. It then proves that the replicated query lies
inside the fetched half-open box
$[\mathbf{N}^{\leftarrow},\mathbf{N}^{\rightarrow})$ coordinate-wise, which
certifies that the fetched row is exactly the accepting leaf of that tree on
that sample. Finally, the fetched leaf scores are aggregated back along the
tree axis to obtain the class scores. The remaining forest objects are not used
directly in this fetch-and-containment check, but they are needed to certify
that the committed rows do form a valid hyperbox forest and to connect
inference with the later training proof.

\subsection{Training Proof Pipeline}

The training proof reuses the same committed forest and is organized around one
logical statement per boosting round. For every round $t\in[T]$ and class
$k\in[K]$, the prover exposes a global leaf-assignment vector
\[
\boldsymbol{\ell}^{(t,k)} \in
\{N_{\mathrm{int}},\ldots,N_{\mathrm{tot}}-1\}^{M},
\]
where $\boldsymbol{\ell}^{(t,k)}[r]$ is the global leaf index reached by sample
$r$ in the tree trained for class $k$ at round $t$.

A key structural point is that the family
$(\mathbf{S}^{(0)},\mathbf{S}^{(1)},\ldots,\mathbf{S}^{(T)})$ is not committed
round by round. Instead, all score matrices are arranged along one extra round
axis and committed in a single PCS instance; the same shared commitment pattern
is used for the aligned boosting-state tensors derived from them. The proof
then checks all score-update, softmax, and residual identities against that
shared committed object. Consequently, the proof size does not grow with the
number of boosting rounds; increasing $T$ enlarges the witness tensor and the
prover's work, but not the PCS-level proof artifact.

Logically, the proof of training establishes that for every round $t\in[T]$,
all of the following conditions hold:
\begin{enumerate}
  \item the current boosting state is correct:
  $\mathbf{S}^{(t)}\in\mathbb{F}^{M\times K}$ is the score matrix carried into
  round $t$, and
  $\mathbf{G}^{(t)}\in\mathbb{F}^{M\times K}$ is the pseudo-residual matrix
  derived from $\mathbf{S}^{(t)}$ by the softmax relation;
  \item for every class-specific tree $(t,k)$, the vector
  $\boldsymbol{\ell}^{(t,k)}$ records exactly which leaf each sample reaches in
  the committed tree of that round;
  \item the histogram tensors derived from $\boldsymbol{\ell}^{(t,k)}$ and
  $\mathbf{G}^{(t)}[\cdot,k]$ are correct, first on leaf rows and then after
  propagation to internal nodes through the routing matrices; these tensors lie
  in $\mathbb{F}^{N_{\mathrm{tot}}\times d\times B}$ over the node, feature,
  and bin axes;
  \item from those histograms, the protocol derives the candidate split-score
  tensor
  $\Gamma^{(t,k)}\in\mathbb{F}^{N_{\mathrm{int}}\times d\times B}$ for every
  internal-node / feature / threshold triple, and the split encoded by
  $(\mathbf{E},\boldsymbol{\Theta})$ is certified to optimize that tensor;
  \item the same per-leaf aggregates determine a padded leaf-value matrix
  $\widetilde{\mathbf{W}}^{(t)}\in\mathbb{F}^{N_{\mathrm{tot}}\times K}$ for
  the round-$t$ trees, and these leaf values update the boosting state to the
  next score matrix $\mathbf{S}^{(t+1)}\in\mathbb{F}^{M\times K}$.
\end{enumerate}

Together, these conditions imply that the committed objects encode a valid
plaintext GBDT training trajectory across all rounds. The tree-specific novelty
lies in the middle of this statement: instead of proving a sorting procedure,
the protocol proves the statistics from which split search can be reconstructed
algebraically.

\subsection{Objective Rewriting for Split Search}

Consider one node $u$, one feature $j$, and one threshold bin $b$. Let
$I_u \subseteq [M]$ be the samples reaching node $u$, and let
$L_{u,j,b},R_{u,j,b}$ be the induced left and right partitions. If each sample
$r\in I_u$ carries target value $g_r$, then the residual-sum-of-squares
objective is
\[
\operatorname{RSS}(u,j,b)
=
\sum_{r\in L_{u,j,b}}(g_r-\bar{g}_L)^2
+
\sum_{r\in R_{u,j,b}}(g_r-\bar{g}_R)^2.
\]
Expanding the squares gives
\begin{equation}\label{eq:rss-rewrite}
\operatorname{RSS}(u,j,b)
=
\sum_{r\in I_u} g_r^2
-
\left(
\frac{\left(\sum_{r\in L_{u,j,b}} g_r\right)^2}{|L_{u,j,b}|}
+
\frac{\left(\sum_{r\in R_{u,j,b}} g_r\right)^2}{|R_{u,j,b}|}
\right).
\end{equation}
For fixed node $u$, the first term in \eqref{eq:rss-rewrite} is independent of
the candidate split. Therefore minimizing $\operatorname{RSS}(u,j,b)$ is
equivalent to maximizing
\begin{equation}\label{eq:split-objective}
\Phi(u,j,b)
:=
\frac{\left(\sum_{r\in L_{u,j,b}} g_r\right)^2}{|L_{u,j,b}|}
+
\frac{\left(\sum_{r\in R_{u,j,b}} g_r\right)^2}{|R_{u,j,b}|},
\end{equation}
subject to the validity condition that both children are nonempty.

This rewriting explains why histograms are sufficient proof objects. Once the
prover has committed, for every node $u$, feature $j$, and bin $q$, the sample
count and residual sum associated with that bin, the quantities
$|L_{u,j,b}|$, $|R_{u,j,b}|$,
$\sum_{r\in L_{u,j,b}} g_r$, and $\sum_{r\in R_{u,j,b}} g_r$ become prefix- and
total-sum expressions over the bin axis. The split objective is then reduced to
tensor algebra over committed counts and sums,
which becomes much easier to handle.

\subsection{Domain-Lifted and Interleaved Batching}

The construction repeatedly needs to batch claims that are structurally similar
but do not initially share one proof shape. In this subsection, we describe two
general-purpose batching techniques that are used throughout the construction
to achieve succinctness and modularity. Prior proof systems already rely on
vanishing-polynomial constraints and random-linear-combination batching on a
common domain \cite{aurora19,cq++24}. Our claim is narrower: to the best of our
knowledge, the two batching moves isolated here are new as reusable
proof-composition tools in this setting. The first lifts heterogeneous linear
claims to the same largest domain before batching, and the second interleaves
aligned same-form nonlinear bundles into one pointwise check.

\paragraph{Domain-lifted batch.}
Most linear constraints in the paper are first expressed over some evaluation
domain $\mathbb{H}_n$. After moving all terms to one side, we obtain a virtual polynomial $F(X)$, and the constraint is
reduced to a zero-check over that domain, hence to a divisibility claim of the
form
\[
F(X)=Q(X)\nu_n(X),
\]
where $\nu_n(X)=X^n-1$ is the vanishing polynomial of $\mathbb{H}_n$.

The basic batching problem is therefore the following: suppose we have two such
claims, one over $\mathbb{H}_{n_1}$ and one over $\mathbb{H}_{n_2}$, with
$n_1<n_2$. Since our domain sizes are powers of two, the smaller domain is
nested inside the larger one. We can therefore lift the smaller claim to
$\mathbb{H}_{n_2}$ by multiplying its residual by
$\nu_{n_2}(X)/\nu_{n_1}(X)$:
\[
\widehat{F}_1(X):=F_1(X)\frac{\nu_{n_2}(X)}{\nu_{n_1}(X)}.
\]
Now both claims live over the same ambient domain $\mathbb{H}_{n_2}$, and the
verifier can batch them as
\[
\widehat{F}_1(X)+\eta F_2(X)=Q(X)\nu_{n_2}(X)
\]
for a random challenge $\eta$.

We call this \emph{domain-lifting batching}. It lets us merge linear claims
that are initially written over different domains by first rewriting the
smaller one over the larger domain and then applying an ordinary same-domain
batch.

\paragraph{Interleaving batching.}
The second batching problem concerns same-form nonlinear checks that are
\emph{local}: validity at one position depends only on the witness values at
that position, not on any other entry of the vector. Typical examples in our
setting include lookup-style checks against one common table and range checks
against one public range table. If each local instance is proved per PCS, the
proof size
scales with the number of instances even though the check itself is always the
same.

The basic case is two vectors of the same length. Suppose
$\mathbf{f},\mathbf{g}\in\mathbb{F}^n$ must both satisfy the same local
constraint, say $\varphi(x)=0$ entrywise. Then we can interleave them into one
larger virtual vector
\[
(\mathbf{f}[0],\mathbf{g}[0],\mathbf{f}[1],\mathbf{g}[1],\ldots),
\]
and check the same local constraint once on that merged vector. In polynomial
form, if $f(X)$ and $g(X)$ are the encodings of $\mathbf{f}$ and $\mathbf{g}$
over $\mathbb{H}_n$, the merged object is
\[
\operatorname{Int}_n(f,g)(X)
:=
\operatorname{even}_n(X)\,f(X)
+ \operatorname{odd}_n(X)\,g(\omega_{2n}^{-1}X),
\]
where
\[
\operatorname{even}_n(X):=\frac{X^n+1}{2},
\qquad
\operatorname{odd}_n(X):=\frac{1-X^n}{2}.
\]
These public sparse selectors place $\mathbf{f}$ on the even positions and
$\mathbf{g}$ on the odd positions, so every position of the larger vector still
corresponds to one concrete original entry.

If the two vectors have different lengths, we first lift the smaller one to the
larger domain by the substitution $X\mapsto X^r$, exactly as in \cref{lem:row-replication}.
The important point is that the merged object is still virtual: it is assembled
from existing commitments, public sparse masks, and domain-lift substitutions,
so no new base commitment is required.

\section{Proof-Batching Tools}\label{sec:tools}

This section presents the two batching tools used repeatedly in the protocol.
The first handles heterogeneous linear identities whose natural domains
differ. The second merges many same-form local nonlinear checks into one
larger virtual instance. Together they cover the two batching patterns that
recur throughout the rest of the construction.
Standard divisibility-based zero checks and same-domain batching already appear
in prior proof systems such as Aurora and CQ++ \cite{aurora19,cq++24}. Our
contribution here is to isolate two reusable tools: domain-lifting batching for
heterogeneous linear claims, and interleaving batching for aligned same-form
nonlinear claims.

\subsection{Domain-Lifting Batching}

As explained earlier, most linear constraints in our construction are first
expressed as zero-checks over some evaluation domain $\mathbb{H}_n$. The
following theorem shows how to batch many such claims together even when their
natural domains differ. The key idea is to lift each claim from its original
domain to a common larger domain, so the lifted claim is still a zero-check
but now over the same domain as the others. After lifting, we can apply an
ordinary same-domain batch to combine all claims into one.

\begin{theorem}[Domain-lifting batching]
\label{thm:domain-lifting-batching}
Let
\[
2 \leq n_1 \leq n_2 \leq \cdots \leq n_m = N
\]
be powers of two. For each $i \in [m]$, let $F_i(X)$ be a virtual polynomial
and suppose the target claim is
\[
F_i(X)=Q_i(X)\nu_{n_i}(X)
\]
for some polynomial $Q_i(X)$. Define the lifted polynomial
\[
\widehat{F}_i(X):=
F_i(X)\frac{\nu_N(X)}{\nu_{n_i}(X)}.
\]
Then:
\begin{enumerate}
  \item $F_i(X)=Q_i(X)\nu_{n_i}(X)$ holds if and only if
  $\widehat{F}_i(X)$ is divisible by $\nu_N(X)$.
  \item For a verifier challenge $\eta$, all $m$ claims can be combined into
  the single identity
  \begin{equation}\label{eq:domain-lifting-batching}
  \sum_{i=1}^{m}\eta^{i-1}\widehat{F}_i(X)=Q(X)\nu_N(X).
  \end{equation}
  \item If at least one original claim is false, then
  \eqref{eq:domain-lifting-batching} holds for at most $m-1$ values of $\eta$, i.e. the soundness error is at most $(m-1)/|\mathbb{F}|$.
\end{enumerate}
\end{theorem}

\begin{proof}
The first item is immediate from
\[
F_i(X)=Q_i(X)\nu_{n_i}(X)
\iff
F_i(X)\frac{\nu_N(X)}{\nu_{n_i}(X)} = Q_i(X)\nu_N(X).
\]
For soundness, once the lifted residuals are reduced modulo $\nu_N(X)$, the
left-hand side of \eqref{eq:domain-lifting-batching} becomes a polynomial in $\eta$
of degree at most $m-1$. If at least one original claim is false, this
polynomial is nonzero, so a uniformly random $\eta$ causes accidental
acceptance with probability at most $(m-1)/|\mathbb{F}|$.
\end{proof}

Theorem~\ref{thm:domain-lifting-batching} is the protocol's main tool for batching
heterogeneous linear identities. Later we apply it to routing equations,
padding-induced tensor equalities, and other zero-check-style relations whose
natural domains are not identical.

\subsection{Interleaving Batching}

The intended setting for interleaving batching is that many checks share the
same local rule and the same public parameters, but arrive as separate 
PCS and witness on possibly different power-of-two domains. We first record the alignment lemma
used to place smaller objects on larger domains.

\begin{lemma}[Row replication]
\label{lem:row-replication}
Let $\mathbf{v}\in\mathbb{F}^m$ and let $v(X)$ be its encoding over
$\mathbb{H}_m$. Let $\mathbf{V}\in\mathbb{F}^{r\times m}$ be the matrix whose
every row equals $\mathbf{v}$. Then the encoding polynomial of $\mathbf{V}$ over
$\mathbb{H}_{rm}$ is
\[
V(X)=v(X^r).
\]
\end{lemma}

\begin{proof}
For every $i\in[r]$ and $j\in[m]$,
\[
V(\omega_{rm}^{im+j})
= \mathbf{V}[i,j]
= \mathbf{v}[j]
= v(\omega_m^j)
= v((\omega_{rm}^{im+j})^r),
\]
where the last equality uses $\omega_{rm}^r=\omega_m$. Since both sides have
degree less than $rm$ and agree on $\mathbb{H}_{rm}$, they are identical.
\end{proof}

\begin{definition}[Pairwise interleaving]
\label{def:pairwise-interleaving}
Let $f(X)$ and $g(X)$ be the encoding polynomials of two length-$n$ vectors
over $\mathbb{H}_n$, where $n$ is a power of two. Define
\[
\operatorname{even}_n(X):=\frac{X^n+1}{2},
\qquad
\operatorname{odd}_n(X):=\frac{1-X^n}{2}.
\]
Their evaluations on $\mathbb{H}_{2n}$ are the alternating selectors
$(1,0,1,0,\ldots)$ and $(0,1,0,1,\ldots)$. The pairwise interleaving of
$f(X)$ and $g(X)$ is defined as the virtual polynomial
\[
\operatorname{Int}_n(f,g)(X)
:=
\operatorname{even}_n(X)\,f(X)
+
\operatorname{odd}_n(X)\,g(\omega_{2n}^{-1}X).
\]
Its evaluation vector over $\mathbb{H}_{2n}$ is
\[
(f[0],g[0],f[1],g[1],\ldots,f[n-1],g[n-1]).
\]
\end{definition}

\begin{theorem}[Interleaving batching]
\label{thm:interleaving-batch}
For each $i\in[m]$, let $u_i(X)$ denote the polynomial encoding a vector of
length $n_i$ over $\mathbb{H}_{n_i}$, where each $n_i$ is a power of two. For
some (non-linear) constraint $\varphi$, suppose the target claim is that all
$m$ pointwise relations
\[
\forall i \in [m],\ \forall a\in\mathbb{H}_{n_i}\quad
\varphi\!\left(
u_i(a)
\right)=0.
\]
There exists an algorithm that combines these $m$ claims into one
check over a single virtual polynomial $\widehat{u}(X)$ over domain
$\mathbb{H}_{\widehat{n}}$, using only domain lifts and pairwise interleaving
and introducing no new PCS commitment. Using two scans over power-of-two
domain sizes, it runs in $O(m+\log \hat{n})$ time. This algorithm guarantees:
\begin{enumerate}
  \item all original pointwise claims hold if and only if
  \[
  \forall z\in\mathbb{H}_{\widehat{n}}\quad
  \varphi\!\left(
  \widehat{u}(z)
  \right)=0;
  \]
  \item the final domain size satisfies
  \[
  \widehat{n}
  =
  2^{\left\lceil \log_2(n_1+\cdots+n_m)\right\rceil}.
  \]
  \item any opening claim for $\widehat{u}(X)$ is reduced to one opening for each
  original $u_i(X)$, and the verifier can check the claim in
  $O(m\log \widehat{n})$ time.
\end{enumerate}
\end{theorem}

\begin{proof}
Start from the active multiset $\{(u_i(X),n_i)\}_{i\in[m]}$.
We first store all pairs with same domain size in the same bucket.
Then we perform two scans over the buckets.
In the first scan, process the buckets from small to large domain size and
merge equal-sized pairs whenever a bucket contains at least two active pairs.
Each such merge uses pairwise interleaving as in
\cref{def:pairwise-interleaving}. After this pass, every bucket contains at
most one active pair, and the total size $\sum_i n_i$ is unchanged.

In the second scan, repeatedly take the two smallest active pairs. If their
sizes differ, lift the smaller one to the larger domain using
$X\mapsto X^r$ with $r=n_j/n_i$; then interleave the pair. Repeat
until one active pair $(\widehat{u}(X),\widehat{n})$ remains.

Because the domain sizes are powers of two, the two scans touch only
$O(\log \hat{n})$ buckets, and each merge or lift removes one active pair, so the
implementation runs in $O(m+\log \hat{n})$ time.

For property~(1), lifting is exactly the setting of \cref{lem:row-replication}:
if $r=n_j/n_i$, then $X\mapsto X^r$ repeats the entries of $u_i(X)$ on the
larger domain $\mathbb{H}_{n_j}$. So the pointwise relation $\varphi(\cdot)=0$
holds before lifting if and only if it holds after lifting. Pairwise
interleaving preserves the same relation for the same reason,
since all entries are preserved.
Repeating these two facts over the two scans proves property~(1).

For property~(2), note that the first scan preserves the total size
$S:=n_1+\cdots+n_m$. If we denote the
sizes after the first scan by $n_1'<\cdots<n_s'$, then
\[
\sum_{i=1}^m n_i = \sum_{i=1}^s n_s'\implies n_s' = 2^{\left\lfloor \log_2(n_1+\cdots+n_m)\right\rfloor}.
\]
If $n_s' = 2^{\log_2(n_1+\cdots+n_m)}$,
there $s=1$ and there is no second scan, and the property~(2) holds.
Otherwise,
the second scan pushes the size to $2 n_s'$,
and thus
\[
\widehat{n}= 2 n_s' = 2^{\left\lceil \log_2(n_1+\cdots+n_m)\right\rceil}.
\]

For property~(3), note that for interleaving,
the verifier can compute the sparse selectors by itself.
And for row-replication $u(X)\mapsto u(X^r) = v(X)$, the opening of $v(x)$ for any $x\in \mathbb{F}$ is trivially reduced to the opening of $u(x^r)$.
Hence, the verifier can check any opening claim for $\widehat{u}(X)$ by checking one opening claim for each original $u_i(X)$,
which incurs $O(m\log \widehat{n})$ time to compute the selectors and check the claims.
\end{proof}

\subsection{Explicit Lookup Batching Example}

As a concrete application for interleaving batching, suppose we have several lookup families indexed by
$i\in[s]$. For each family, there is a public table $\mathbf{t}_i$ and a
collection of vector queries $\mathbf{f}_{i,1},\ldots,\mathbf{f}_{i,m_i}$ such
that every query satisfies $\mathbf{f}_{i,j}\subseteq \mathbf{t}_i$.

We batch them in three steps. First, for each fixed $i$, we interleave the
queries $\mathbf{f}_{i,1},\ldots,\mathbf{f}_{i,m_i}$ into one family-level
query vector $\mathbf{f}_i$. Second, we shift each table by a random challenge
$\alpha$ and form $\widetilde{\mathbf{t}}_i:=\alpha^i+\mathbf{t}_i$, so that
different tables are separated with high probability. We then interleave the
shifted tables into one combined table vector $\mathbf{t}$. Third, we
interleave the family-level query vectors $\mathbf{f}_1,\ldots,\mathbf{f}_s$
into one global query vector $\mathbf{f}$.

This yields a single lookup claim
\[
\mathbf{f}\subseteq \mathbf{t},
\]
which can be proved with one lookup argument instead of one argument per
family. The benefit is that the prover amortizes the lookup proof across all
tables while keeping the query structure explicit. The random shifts prevent
cross-table collisions, and the interleaving steps preserve the original
within-family lookup relations.

A second example is fixed-point quantization. Suppose several checks all prove
relations of the form
\[
a_i = \left\lfloor \frac{b_i}{2^q} \right\rfloor
\]
for the same public divisor $2^q$ and the same rounding rule. These are again
same-form local checks: after aligning the domains of the witness vectors, we
interleave them and obtain one larger quantization check. This is the pattern
behind the fixed-point gain, leaf-value, and normalization relations that
appear later in the construction.

Theorem~\ref{thm:interleaving-batch} is useful when the protocol produces many
small nonlinear checks of the same form, and lookup batching across different
tables is one of the main examples in our construction.

\section{Evaluation}

The empirical evaluation is intended to measure the cost of the primitive layer
and the end-to-end protocol separately. At the primitive level, the most
relevant metrics are the prover time, verifier time, and proof size of domain
lifting, replication, tensor sums, aggregation, and histogram construction. At
the application level, the evaluation reports the cost of batched inference and
end-to-end training proofs as functions of the number of trees, the maximum
depth, the feature dimension, the number of bins, and the batch size.

\subsection{Experimental Methodology}

The experiments compare the proposed protocol against the most relevant
baselines suggested by the literature and by the proof architecture itself.
Natural baselines include direct reuse of decision-tree-oriented lookup systems
\cite{cq++24}, data-parallel zkSNARK approaches to tree proving
\cite{sparrow24}, and recent training-oriented baselines such as
\cite{zkboost26}. Of particular interest is whether the histogram and
aggregation layer yields the intended improvement over proof strategies that
explicitly sort, compare, or RAM-simulate candidate splits inside a generic
backend.

\section{Related Work}

Organize prior work by topic instead of paper-by-paper summary when possible:
\begin{itemize}
  \item verifiable machine learning,
  \item zero-knowledge inference,
  \item decision-tree and GBDT verification,
  \item privacy-preserving tree-based learning.
\end{itemize}

You can start this section with a short framing paragraph such as:
recent work on verifiable machine learning and tree-based models provides
useful baselines for positioning a ZK-GBDT system
\cite{placeholder-zkml,placeholder-gbdt}.

\section{Conclusion}\label{sec:conclusion}

This paper formulates a proof architecture for quantized GBDT in which the
central training difficulty, namely optimal split selection, is handled through
a dedicated primitive layer rather than through a monolithic generic circuit.
By committing the training trace as structured tensors, \textsc{Terrae} proves
routing, histogram construction, split optimality, leaf-value assignment, and
boosting-state updates directly over polynomial commitments. Its main technical
tools, domain-lifting batching and interleaving batching, reduce many
heterogeneous algebraic and non-native checks to a small number of batched
claims without introducing extra auxiliary commitments. The histogram
aggregation argument further avoids a generic RAM proof for the dominant
sample-to-bucket update pattern in tree training. Together, these components
give a specialized zero-knowledge proof system for GBDT training and inference,
and the evaluation shows that exploiting this structure substantially reduces
prover cost relative to a circuit-style baseline.


\bibliographystyle{plain}
\bibliography{bibliography/references}

\end{document}
